%!TEX root = TheoretischePhysikC.tex


\section{Vorlesung: Zerfließen des Wellenpakets}
Der Ausdruck ,,Wellenpaket`` wird für räumlich begrenzte Funktionen verwendet, wie etwa $\Psi\in L^2(\mathbb{R}^3)$, nicht für ebene Wellen. Vom Standpunkt der statistischen Interpretation gibt es also nur solche ,,Pakete``, und die Diskussion ebener Wellen ist bestenfalls eine etwas verkürzte Art, über die Fourier-Transformation zu sprechen. In der Frühzeit der Quantenmechanik galten gut lokalisierte Wellenpakete außerdem als eine Art Ersatz für Teilchen. Gemessen
daran war eine Eigenschaft der Wellenpakete etwas störend: sie ,,flossen auseinander``,  was für Teilchen ja wirklich ein ungebührliches Betragen ist. Das Zerfließen des Wellenpakets galt danach als typisch quantenmechanischer Effekt. Wenn man allerdings $\Psi$ konsequent statistisch deutet, ist nichts Paradoxes an diesem Vorgang: warum soll schließlich eine Wahrscheinlichkeitsverteilung
nicht mit der Zeit breiter werden? Wir werden hier sehen, dass das Zerfließen der Wellenpakete gerade ein klassischer Effekt ist: er wird asymptotisch genau durch die Impulsverteilung bestimmt, wie bei einem Schwarm klassischer Teilchen.

\begin{greybox}[Warum reale Teilchen trotzdem nicht zerfließen]
	Allein ist ein Teilchen selten. Materie besteht aus geladenen Bestandteilen, und das elektromagnetische Feld fragt sie unablässig ,,bist du da?`` -- das Teilchen antwortet, indem es ein wenig Feld erzeugt, und diese pausenlose Ortsmessung lokalisiert es immer wieder neu. Stichwort: \emph{Dekohärenz}.
\end{greybox}

Wir betrachten also jetzt Teilchen in $\mathbb{R}^3$ mit Hamiltonoperator
\begin{equation}
	H = \frac{\widehat{p}_\text{ges.}^2}{2m} = \frac{\widehat{p}_{x}^2}{2m}+\frac{\widehat{p}_{y}^2}{2m}+\frac{\widehat{p}_{z}^2}{2m}.
\end{equation}
In diesem Fall entspricht wie üblich in der QM $H$ der Energie eines freien Teilchens. 

Wir interessieren uns für die Wahrscheinlichkeitsverteilung des Ortes zu verschiedenen Zeiten, besonders zu großer Zeit, das heißt, für die Ergebnisse von
Ortsmessungen, wenn zwischen Präparation und Messung eine große Zeit $t$ verstrichen ist. Wir haben zwei äquivalente Alternativen, um dieses Problem anzugehen: das Heisenbergbild und das Schrödingerbild. Im Schrödingerbild würden wir die Schrödingergleichung 
\begin{equation}
	i\hbar \frac{\partial}{\partial t}\Psi = \widehat{H}\Psi
\end{equation}
lösen, und dann das Ergebnis als 
\begin{equation}
	p_t(x) \equiv |\Psi(x,t)|^2
\end{equation}
ablesen -- eine Wolke im Raum, Detektorklick für Detektorklick aufgesammelt. Der Unterschied der Bilder ist reine Buchhaltung: dort bewegt sich der Zustand, hier die Messung, und wegen $\tr(\rho_t\widehat{F}) = \tr(\rho\,\widehat{F}(t))$ liefern beide dieselben Zahlen. Wir werden die Rechnung im Heisenbergbild durchführen, d.h. wir halten $\Psi$ (und dabei
die Dichtematix $\rho$) fest und berechnen die Operatoren
\begin{equation}
	\widehat{x}(t) \equiv e^{\frac{it}{\hbar} \widehat{H}}\widehat{x} e^{-\frac{it}{\hbar} \widehat{H}}, \quad \widehat{y}(t) \equiv e^{\frac{it}{\hbar} \widehat{H}}\widehat{y} e^{-\frac{it}{\hbar} \widehat{H}} \quad \text{und} \quad \widehat{y}(t) \equiv e^{\frac{it}{\hbar} \widehat{H}}\widehat{y} e^{-\frac{it}{\hbar} \widehat{H}}.
\end{equation}

Als nächstes suchen wir passende ,,Koordinaten``, in denen wir unsere Rechnung duchführen können. Wegen des Auftretens der Exponentialfunktion ist es naheliegend, eine Darstellug des Hilbertraums zu wählen, in der $H$ ein Multiplikationsoperator ist, also die Impulsdarstellung:
\begin{equation}
	\Phi(\mathbf{p})\equiv \Phi(p_x,p_y,p_z) \equiv \frac{1}{\sqrt{2\pi\hbar}} \iiint_{-\infty}^\infty \Psi(x)e^{-\frac{i}{\hbar}(p_x x+ p_y y+ p_z z)}\,dxdydz.	
\end{equation} 
Zu jedem $\mathbf{p}$ gehört dabei die ebene Welle $|\Psi_{\mathbf{p}}\rangle$ aus Vorlesung 8, hier so normiert, dass $\langle\Psi_{\mathbf{p}}|\Psi_{\mathbf{p}'}\rangle = \delta^{(3)}(\mathbf{p}-\mathbf{p}')$ gilt:
\begin{equation}
	\langle \mathbf{x}|\Psi_{\mathbf{p}}\rangle = \frac{1}{(2\pi\hbar)^{3/2}}e^{\frac{i}{\hbar}(p_x x + p_y y + p_z z)},
\end{equation}
und diese uneigentlichen Kets liefern eine Zerlegung der Eins:
\begin{equation}\label{eq:eins-impuls}
	\mathbb{I} = \iiint_{-\infty}^\infty |\Psi_{\mathbf{p}}\rangle\langle\Psi_{\mathbf{p}}|\,dp_xdp_ydp_z.
\end{equation}
Die Koeffizienten $\langle\Psi_{\mathbf{p}}|\Psi\rangle$ in dieser Basis sind gerade $\Phi(\mathbf{p})$.\footnote{Der Vorfaktor oben in $\Phi(\mathbf{p})$ ist der eindimensionale; dreidimensional gehört dorthin $(2\pi\hbar)^{-3/2}$.} Solches Einschieben einer Eins ist unser wichtigstes Rechenwerkzeug; wir brauchen es gleich wieder.

Gewählt ist diese Basis, weil die ebenen Wellen Eigenvektoren von $\widehat{H}$ sind,
\begin{equation}\label{eq:heigen}
	\widehat{H}|\Psi_{\mathbf{p}}\rangle = \frac{|\mathbf{p}|^2}{2m}|\Psi_{\mathbf{p}}\rangle,
\end{equation}
so dass in der Potenzreihe jedes $\widehat{H}$ zu einer Zahl wird:
\begin{equation}
	e^{-\frac{it}{\hbar}\widehat{H}}|\Psi_{\mathbf{p}}\rangle = \sum_{n=0}^{\infty}\frac{1}{n!}\left(-\frac{it}{\hbar}\right)^n \widehat{H}^n|\Psi_{\mathbf{p}}\rangle = e^{-\frac{it}{\hbar}\frac{|\mathbf{p}|^2}{2m}}|\Psi_{\mathbf{p}}\rangle.
\end{equation}
Im Exponenten stand ein Operator, jetzt eine Zahl -- daher unsere Besessenheit von Eigenvektoren. Erlaubt ist das \emph{nur} auf einem Eigenvektor: Eine Matrix durch eine Zahl zu ersetzen, das muss man sich verdienen.

Wir brauchen gleich außerdem, dass $\widehat{x}$ in der Impulsdarstellung als $i\hbar\frac{\partial}{\partial p_x}$ wirkt. Eindimensional schiebt man dazu die Ortszerlegung der Eins ein,
\begin{equation}
	\langle \Psi_p|\widehat{x}|\Psi\rangle = \int_{-\infty}^\infty \langle \Psi_p|x\rangle\, x\, \Psi(x)\,dx = \frac{1}{\sqrt{2\pi\hbar}}\int_{-\infty}^\infty x\, e^{-\frac{i}{\hbar}px}\,\Psi(x)\,dx,
\end{equation}
erzeugt den Faktor $x$ als Ableitung nach $p$,
\begin{equation}
	i\hbar\frac{\partial}{\partial p}e^{-\frac{i}{\hbar}px} = i\hbar\left(-\frac{ix}{\hbar}\right)e^{-\frac{i}{\hbar}px} = x\,e^{-\frac{i}{\hbar}px},
\end{equation}
und zieht diese vor das Integral:
\begin{equation}\label{eq:xinp}
	\langle \Psi_p|\widehat{x}|\Psi\rangle = i\hbar\frac{\partial}{\partial p}\Phi(p).
\end{equation}
Für $\widehat{y}$ und $\widehat{z}$ läuft es genauso -- eine schöne Symmetrie zur Ortsdarstellung, wo umgekehrt $\widehat{p}$ als $-i\hbar\frac{\partial}{\partial x}$ wirkt: dieselbe Struktur mit vertauschten Rollen, bis auf das Vorzeichen.

Dann ist
\begin{equation}\label{eq:timedepx}
	\begin{split}
		(\widehat{x}(t)\Phi)(\mathbf{p}) &= e^{\frac{it}{\hbar} \frac{|\mathbf{p}|^2}{2m}}\left(\widehat{x}(t)e^{-\frac{it}{\hbar}\widehat{H}}\Phi\right)(\mathbf{p}) \\
		&= e^{\frac{it}{\hbar} \frac{|\mathbf{p}|^2}{2m}}\left(i\hbar \frac{\partial}{\partial p_x}\right)\left(e^{-\frac{it}{\hbar}\frac{|\mathbf{p}|^2}{2m}}\Phi(\mathbf{p})\right) \\
		&= t \frac{p_x}{m}\Phi(\mathbf{p}) + i\hbar \frac{\partial}{\partial p_x} \Phi(\mathbf{p}), \\ 
	\end{split}
\end{equation}
also
\begin{equation}
	\widehat{x}(t) = t \frac{\widehat{p}_x}{m} + \widehat{x}(0),
\end{equation}
wobei $\frac{\widehat{p}_x}{m}$ den \emph{Geschwindigkeitsoperator} (in der $x$-Richtung) darstellt. Es gelten ähnliche Ausdrücke für $\widehat{y}(t)$ und $\widehat{z}(t)$. Den Ausdruck $\frac{p_\mu}{m} = \frac{\partial}{\partial p_\mu} \frac{|\mathbf{p}|^2}{2m}$, $\mu = x,y,z$, bezeichnet man auch als \emph{Gruppengeschwindigkeit}, weil er die Geschwindigkeit von \emph{Gruppen}, also Wellenpaketen beschreibt. Die Rechnung ist aufwendig; merken sollte man sich ihr Ergebnis, nicht die Schrittfolge.

Um die Zeitabhängigkeit des Erwartungswertes von $\widehat{x}(t)$ zu bekommen, braucht man nun nur noch auf beiden Seiten von (\ref{eq:timedepx}) die Spur mit $\rho =  |\Psi\rangle\langle\Psi|$ zu nehmen. Wir erhalten dabei
\begin{equation}
	\begin{split}
		\langle\widehat{x}\rangle(t) &= \tr(\rho \widehat{x}(t)) \\
		&= \tr\left(\rho \left(t \frac{\widehat{p}_x}{m} + \widehat{x}(0)\right)\right)\\
		&= t \frac{\langle\widehat{p}_x \rangle(0)}{m} + \langle\widehat{x}\rangle(0).
	\end{split}
\end{equation}
Das ist genau die newtonsche Regel ,,Ort gleich Anfangsort plus Geschwindigkeit mal Zeit``.

\begin{intuition}
	Dieselbe Formel, andere Physik: Klassisch \emph{hat} ein Teilchen eine Position, Erwartungswerte braucht man gar nicht. Hier läuft nur der Mittelwert klassisch, während die Breite wächst -- eine klassische Wolke mit scharfem Ort und Impuls zerfließt nicht, eine quantenmechanische muss es. Klassische Mechanik ist der Grenzfall $\hbar\to 0$.
\end{intuition}

Die Operatorgleichung enthält aber viel mehr Information als nur das: Die Diagonalisierung der Operatoren $\widehat{x}(t)$ enthält ja alle Information über die Wahrscheinlichkeitsverteilungen des Ortes zur Zeit $t$ in beliebigen Zuständen. Das Aussehen der Ortswahrscheinlichkeitsverteilung zu großen Zeiten ist also völlig von der gemeinsamen Verteilung der Operatoren $\frac{\widehat{p}_\mu}{m}$ bestimmt. Am Verhalten von $\widehat{x}$, $\widehat{y}$ und $\widehat{z}$  ist also überhaupt kein Unterschied zwischen klassischen und quantenmechanischen Teilchen zu sehen.

% fig_vl10_zerfliessen.tex — Zerfliessen des Wellenpakets: the packet
% drifts uniformly while it spreads, and its width settles onto exactly
% the straight line a swarm of classical particles with the same
% momentum spread would follow.
%
% GEOMETRY FROM THE FORMULA.  Everything here follows from the operator
% identity this lecture derives,
%
%       x(t) = x(0) + (t/m) p ,
%
% by taking mean and variance in a state with no initial x-p
% correlation:
%
%       <x>(t)  = <x>(0) + <p> t/m ,
%       Dx(t)^2 = Dx(0)^2 + (Dp t/m)^2 .
%
% With tau = m Dx(0)/Dp — and for a minimum-uncertainty packet,
% Dx(0) Dp = hbar/2, hence tau = 2 m Dx(0)^2/hbar — the second line is
%
%       Dx(t) = Dx(0) sqrt(1 + (t/tau)^2),
%
% the hyperbola of panel (b).  Its asymptote Dp t/m is precisely how a
% classical swarm with momentum spread Dp comes apart, so the claim of
% the lecture — that the spreading is asymptotically fixed by the
% momentum distribution alone — is the statement that the two curves in
% (b) merge.  Panel (a) plots the Gaussian carrying those two moments,
%
%       |Psi(x,t)|^2 = exp(-(x-<x>(t))^2/(2 Dx(t)^2)) / (sqrt(2 pi) Dx(t)),
%
% so every centre, width and height is forced by the formula rather than
% drawn by eye.  Lengths in units of Dx(0), times in units of tau, drift
% fixed by <p> = 3 Dp: the centres then sit at 0, 3, 6, 9 and the widths
% at 1, sqrt2, sqrt5, sqrt10.
%
% DESIGN DECISIONS
%   * colour encodes TIME and nothing else — entries 1, 2, 4, 5 of the
%     ordered `tpclevels` ramp (navy, blue, orange, red), the same four
%     in both panels, so a curve in (a) can be found again as a dot in
%     (b).  Entry 3 (munsell) is skipped on purpose: next to cgblue it
%     is the same teal to the eye, and here the two would be neighbours
%     in time rather than well separated as in fig_vl11;
%   * the width curve in (b) is green because it is a different
%     QUANTITY, not a fifth time; its formula is set in the same green,
%     which spares the panel a legend;
%   * grey is reference furniture only, and every grey annotation is
%     set in grey too: the verticals at <x>(t), whose EQUAL SPACING is
%     uniform motion; the dashed envelope through the peaks, which
%     falls as 1/Dx(t) exactly as fast as the widths grow, so the
%     conserved area is shown rather than asserted, and whose label is
%     set in the empty strip immediately above it rather than in the
%     right margin, where it crowded the t=3tau label; the initial width
%     and the classical asymptote in (b);
%   * time labels sit to the RIGHT of each peak rather than above it,
%     which is the only side on which the descending envelope leaves
%     room;
%   * panel (b) uses `scale only axis` with a 4.2 x 5.4 cm box on ranges
%     4.2 x 5.4, i.e. 1 cm per unit on both axes.  The asymptote is then
%     a true 45-degree line and its label can be set along it.  The y
%     range runs to 5.4 although the curve reaches only sqrt(1+4.2^2) =
%     4.32 inside the frame: that headroom is where the three annotation
%     lines sit, so none of them is written over the curve.  The widest
%     of them, the green formula, is 3.8 cm in the fonts of the notes —
%     wider than in the cmfonts preview, which is why the panel is sized
%     against the former;
%   * both panels are 5.4 cm tall so that titles and x axes line up;
%   * the orders of magnitude for tau — 10^-16 s for an electron
%     localised to an angstrom, 10^13 s for a microgram grain — are
%     NOT in the panel.  They were, at its foot, and in the fonts of the
%     notes (wider than the cmfonts preview) the grey asymptote ran
%     straight through them; a 4.2 cm panel crossed by two reference
%     lines has no free region 3.4 cm wide, so the numbers now stand in
%     the sentence below the figure, where the definition of tau in (b)
%     hands them straight over.

\begin{center}
\begin{tikzpicture}

% ─────────────────────────────────────────────────────────────
%  (a)  |Psi(x,t)|^2 at four times
% ─────────────────────────────────────────────────────────────
\begin{axis}[
    name=packets,
    tpcplot,
    scale only axis, width=6.4cm, height=5.4cm,
    grid=none,
    xmin=-3.2, xmax=16.8,
    ymin=0,    ymax=0.47,
    xtick={0,3,6,9,12,15},
    ytick={0,0.1,0.2,0.3,0.4},
    xlabel={$x/\Delta x(0)$},
    ylabel={$\Delta x(0)\,|\Psi(x,t)|^2$},
    every axis x label/.style={at={(axis description cs:0.5,-0.15)},
                               anchor=north, font=\small},
    every axis y label/.style={at={(axis description cs:-0.16,0.5)},
                               rotate=90, anchor=south, font=\small},
    title={(a) \tpcifenglish{The packet spreads}{Das Paket zerfließt}},
    samples=401,
]

  % Centres <x>(t) = 3 t/tau, each rule drawn only up to the peak it
  % belongs to.  Their equal spacing is the newtonian rule.
  \draw[levellinegrey] (axis cs:0,0) -- (axis cs:0,0.3989);
  \draw[levellinegrey] (axis cs:3,0) -- (axis cs:3,0.2821);
  \draw[levellinegrey] (axis cs:6,0) -- (axis cs:6,0.1784);
  \draw[levellinegrey] (axis cs:9,0) -- (axis cs:9,0.1262);

  % Envelope of the peaks: 1/(sqrt(2 pi) Dx(t)) with t/tau = s/3.  It
  % passes exactly through the four maxima.
  \addplot[envelope, domain=0:12.6, samples=140]
      {1/(sqrt(2*pi)*sqrt(1+x^2/9))};

  % The packets.  Widths sqrt(1+(t/tau)^2) = 1, sqrt2, sqrt5, sqrt10.
  % Each runs across the whole panel: a curve stopped short leaves a
  % blunt end lying on the axis, which reads as a cut-off rather than
  % as a tail that has died away.
  \addplot[wavefn, spacecadet,   domain=-3.2:16.8]
      {exp(-(x)^2/2)/sqrt(2*pi)};
  \addplot[wavefn, cgblue,       domain=-3.2:16.8]
      {exp(-(x-3)^2/4)/sqrt(2*pi*2)};
  \addplot[wavefn, birrenorange, domain=-3.2:16.8]
      {exp(-(x-6)^2/10)/sqrt(2*pi*5)};
  \addplot[wavefn, birrenred,    domain=-3.2:16.8]
      {exp(-(x-9)^2/20)/sqrt(2*pi*10)};

  % Direct labels beside the peaks, each in the colour of its curve and
  % set just above the peak height, where the envelope has already
  % dropped away.
  \node[lbl, text=spacecadet,   anchor=south west]
    at (axis cs:0.55,0.410) {$t=0$};
  \node[lbl, text=cgblue,       anchor=south west]
    at (axis cs:3.45,0.293) {$t=\tau$};
  \node[lbl, text=birrenorange, anchor=south west]
    at (axis cs:6.45,0.189) {$t=2\tau$};
  \node[lbl, text=birrenred,    anchor=south west]
    at (axis cs:9.45,0.136) {$t=3\tau$};

  % What the grey verticals say, and what the grey envelope says.
  \node[lbl, text=birrengray, font=\scriptsize, anchor=north east,
        align=right] at (axis cs:16.6,0.468)
    {\tpcifenglish{centres}{Zentren}:
     $\langle\widehat{x}\rangle(t)=\langle\widehat{p}\rangle\,t/m$\\[1pt]
     \tpcifenglish{area}{Fl\"ache}: $\int|\Psi|^2dx=1$};
  \node[lbl, text=birrengray, font=\scriptsize, anchor=south west]
    at (axis cs:4.7,0.232)
    {\tpcifenglish{peak height}{H\"ohe} $\propto 1/\Delta x(t)$};
\end{axis}

% ─────────────────────────────────────────────────────────────
%  (b)  The width: hyperbola against its classical asymptote
% ─────────────────────────────────────────────────────────────
\begin{axis}[
    name=widths,
    tpcplot,
    at={(packets.south east)}, anchor=south west, xshift=17mm,
    scale only axis, width=4.2cm, height=5.4cm,
    xmin=0, xmax=4.2, ymin=0, ymax=5.4,
    xtick={0,1,2,3,4}, ytick={0,1,2,3,4},
    grid=major,
    xlabel={$t/\tau$},
    ylabel={$\Delta x(t)/\Delta x(0)$},
    every axis x label/.style={at={(axis description cs:0.5,-0.15)},
                               anchor=north, font=\small},
    every axis y label/.style={at={(axis description cs:-0.21,0.5)},
                               rotate=90, anchor=south, font=\small},
    title={(b) \tpcifenglish{Width}{Breite} $\Delta x(t)$},
]

  % Reference furniture: the classical asymptote Dp t/m and the width
  % the packet started with.  They cross at t = tau.
  \addplot[levellinegrey, domain=0:4.2, samples=2] {x};
  \draw[levellinegrey] (axis cs:0,1) -- (axis cs:1,1);

  % Dx(t)/Dx(0) = sqrt(1+(t/tau)^2)
  \addplot[birrengreen, line width=1.1pt, domain=0:4.2, samples=140]
      {sqrt(1+x^2)};

  % The four times of panel (a), in their colours.
  \addplot[only marks, mark=*, mark size=2.0pt, spacecadet]
      coordinates {(0,1)};
  \addplot[only marks, mark=*, mark size=2.0pt, cgblue]
      coordinates {(1,1.41421)};
  \addplot[only marks, mark=*, mark size=2.0pt, birrenorange]
      coordinates {(2,2.23607)};
  \addplot[only marks, mark=*, mark size=2.0pt, birrenred]
      coordinates {(3,3.16228)};

  % The curve names itself, in its own colour; the grey lines below it
  % state the asymptotics and fix the time scale.
  \node[text=birrengreen, font=\scriptsize, anchor=north west]
    at (axis cs:0.08,5.32)
    {$\Delta x(t)=\Delta x(0)\sqrt{1+(t/\tau)^{2}}$};
  \node[lbl, text=birrengray, font=\scriptsize, anchor=north west]
    at (axis cs:0.08,4.68)
    {$\longrightarrow\ \Delta p\,t/m$ \tpcifenglish{for}{f\"ur} $t\gg\tau$};
  \node[lbl, text=birrengray, font=\scriptsize, anchor=north west]
    at (axis cs:0.08,4.14) {$\tau=2m\,\Delta x(0)^2/\hbar$};

  % The asymptote, labelled along itself — the panel is drawn at
  % 1 cm per unit on both axes, so the line really does run at 45 deg.
  \node[lbl, text=birrengray, font=\scriptsize, rotate=45,
        anchor=north west, inner sep=1.5pt] at (axis cs:1.65,1.50)
    {\tpcifenglish{classical swarm}{klassischer Schwarm}};
\end{axis}

\end{tikzpicture}
\end{center}


In~(a) rücken die grauen Marken der Zentren in gleichen Abständen vor, während
die Höhe genau so weit fällt, wie die Breite wächst. In~(b) achte man darauf,
wie früh die Kurve auf die klassische Gerade einschwenkt: schon nach wenigen
$\tau$ ist der Unterschied nicht mehr zu sehen. Ob man das Zerfließen je zu
Gesicht bekommt, entscheidet die Zeitskala $\tau$: für ein Elektron mit
$\Delta x(0) = 10^{-10}$\,m ist $\tau \sim 10^{-16}$\,s, für ein Staubkorn von
$1\,\mu$g mit $\Delta x(0) = 10^{-6}$\,m dagegen $\tau \sim 10^{13}$\,s, also
rund $300\,000$ Jahre.

Die interessanteren Unterschiede zwischen klassischer Mechanik und Quantenmechanik (und die wichtigeren Anwendungen der Quantenmechanik) finden sich allerdings bei Teilchen in äußeren Potentialen, oder in Wechselwirkung miteinander.

Der nächste Schritt dorthin ist bescheiden: ein Teilchen in einem \emph{endlich} großen Kasten -- unser erstes Beispiel dafür, was ,,Quanten`` in Quantenmechanik heißt.